\documentclass[10pt]{letter}
\usepackage[a4paper,margin=2.0cm]{geometry}
\usepackage[T1]{fontenc}
\usepackage{lmodern}
\signature{Mohd-Zulhilmi Paiz Ismadi\\Department of Mechanical Engineering, School of Engineering\\Monash University
Malaysia, 47500 Bandar Sunway, Selangor, Malaysia\\mohd.zulhilmi@monash.edu}
\address{}
\date{\today}
\begin{document}
\begin{letter}{Guest Editors, Special Issue ``Integrating Imaging with Personalized Cardiovascular Medicine: Current
Advances and Future Directions''\\IEEE Journal of Biomedical and Health Informatics}
\opening{Dear Guest Editors,}

We submit the manuscript ``Topological Segmentation Error and Boundary-Condition Tuning in Computed Coronary FFR: A
Controlled In Silico Study'' for consideration in the special issue. The call identifies uncertainty propagation,
segmentation errors and robust validation frameworks as central challenges for image-based cardiovascular models. Our
study addresses all three, for computed fractional flow reserve (FFR).

Computed FFR depends on the segmented coronary lumen twice: through the arteries in the model and through outlet
boundary conditions that are derived from those arteries and, increasingly, tuned to match perfusion imaging. We
inserted 150 controlled stenoses into 108 coronary trees from a public CT angiography dataset, applied four segmentation
error types and recomputed FFR under three boundary-condition protocols in two microvascular bed structures, with a
three-dimensional case study of the same mechanism. The main findings are:
\begin{itemize}
\item Topological errors (a missed side branch or a vessel break) changed the treatment decision at 0.80 in 32--33\%
of models with fixed boundary conditions, compared with 6--10\% for caliber errors (a vessel drawn too wide or too narrow) of inter-observer magnitude, which
were close to the variability of repeat invasive measurement.
\item The boundary-condition protocol alone changed how often a topological error altered the decision. On identical anatomy, a vessel break changed the decision in 43\% of models with fixed boundary conditions and in 5\% with re-derived ones.
\item Tuning the boundary conditions to perfusion reduced decision changes but removed the mismatch that marked the
error: 8--19\% of tuned models with a topological error passed a perfusion check stricter than measurement
repeatability while their FFR was wrong by more than 0.05, against 3--4\% for models with correct anatomy.
\end{itemize}

These results bear directly on how image-based models are validated: agreement with perfusion after tuning does not
establish that the computed FFR is correct, and overlap-based segmentation metrics do not capture the error types that
carry the decision risk. The work combines medical image analysis, physiological modeling and decision-level
evaluation, which matches the scope of the Journal and of this special issue.

The manuscript is original, has not been published, and is not under consideration elsewhere. All authors have read
and approved the submission. The authors declare no conflicts of interest, and the work received no external funding.
The study used a public, anonymized dataset (ImageCAS-X, CC BY 4.0). The code and the frozen cohort list are available
from the corresponding author on request. Supplementary Material accompanies the manuscript.

\closing{Yours sincerely,}
\end{letter}
\end{document}
